\begin{table}[htbp]
\centering
\caption{Main pooled specification: border $\times$ entry-regime interaction}
\label{tab:main_pooled}
\footnotesize
\resizebox{\textwidth}{!}{%
\begin{tabular}{lccccccc}
\toprule
& \multicolumn{2}{c}{Border} & \multicolumn{2}{c}{Capped} & \multicolumn{2}{c}{Border $\times$ Capped} & \\
Outcome & Coef. (\%) & $p$ (cl.) & Coef. (\%) & $p$ (cl.) & Coef. (\%) & $p$ (cl.) & $n$ \\
\midrule
Revenue per store & $+0.395$ ($+48.4\%$) & 0.108 & $+1.068$ ($+191.0\%$) & 0.000 & $+0.034$ ($+3.5\%$) & 0.905 & 995 \\
Revenue per capita & $+1.076$ ($+193.3\%$) & 0.000 & $+0.129$ ($+13.8\%$) & 0.165 & $-0.781$ ($-54.2\%$) & 0.008 & 1314 \\
Dispensary density & $+0.697$ ($+100.7\%$) & 0.000 & $-0.217$ ($-19.5\%$) & 0.130 & $-0.818$ ($-55.9\%$) & 0.000 & 995 \\
\bottomrule
\end{tabular}
}
\begin{minipage}{0.98\textwidth}
\vspace{0.3em}
\footnotesize
\textit{Notes:} Pooled across Colorado, Washington, Oregon, Montana, and New Mexico. Coefficients are OLS,
log-linear outcome; \% is $(e^{\hat\beta}-1)\times 100$. $p$ (cl.) is county-clustered. Revenue-per-store and
dispensary-density specs omit state fixed effects because \texttt{entry\_capped} has no within-state variation
in the per-store sample (Oregon's per-store data is a single 2025 cross-section), which would otherwise make
it collinear with the state dummies. Revenue-per-capita includes state FE since Oregon's per-capita panel
spans both its free (2015--2021) and capped (2022--2025) periods, giving genuine within-state identification.
Entry-regime coding: Colorado, Washington, and Montana are capped for their entire panel (local-government
control over licensing -- county opt-out bans in Colorado, opt-in requirements in Montana -- functions as a
binding restriction even where no state-level numerical cap exists); Oregon is free 2015--2021 and capped
2022--2025; New Mexico is free throughout (statutory anti-cap provision, NMSA 1978 \S 26-2C-6(I)).
\end{minipage}
\end{table}
